\section{Wartości własne macierzy.}

W tej sekcji zakładamy, że pracujemy nad ciałem $\mathbb{C}$.


\begin{definition}
	\textit{Wektorem własnym} macierzy $A \in \mathbb{C}^{n \times n}$
		nazywamy niezerowy wektor $v \in \mathbb{C}^{n \times 1}$ taki,
		że dla pewnego $\lambda \in \mathbb{C}$
	\begin{equation}
		Av = \lambda v.
	\end{equation}
	$\lambda$ nazywamy \textit{wartością własną} odpowiadającą wektorowi własnemu $v$.
\end{definition}

Załóżmy,
	że daną mamy macierz $A \in \mathbb{C}^{n \times n}$
	i chcemy wyznaczyć jej wektory i wartości własne.
Jak możemy to zrobić?
Przekształćmy równanie definiujące te pojęcia.
\begin{align}
	Av = \lambda v, \\
	\textbf{0} = \lambda v - Av, \\
	\textbf{0} = \lambda I v - Av, \\
	\textbf{0} = (\lambda I - A) v.
\end{align}
Oznacza to,
	że $v$ jest wektorem własnym macierzy $A$ z odpowiadającą mu wartością własną $\lambda$
	wtedy i tylko wtedy, gdy $(\lambda I - A) v = \textbf{0}$,
	czyli wtedy i tylko wtedy, gdy $v \in \ker(\lambda I - A)$.

Skoro interesują nas tylko takie $v$,
	które są niezerowe,
	to warunek ten może być spełniony wtedy i tylko wtedy,
	gdy jądro $\ker(\lambda I - A)$ jest nietrywialne.
Zachodzi to dokładnie wtedy,
	gdy $A - \lambda I$ jest macierzą nieodwracalną,
	czyli
\begin{equation}\label{eq:eigen-condition}
	\det(\lambda I - A) = 0.
\end{equation}

Aby wyznaczyć wartości własne $A$, musimy więc sprawdzić,
	dla jakich $\lambda \in \mathbb{C}$ spełnione jest równanie \ref{eq:eigen-condition}.
Okazuje się to dość proste, gdyż zawsze jest to równanie wielomianowe (zadanie \ref{ex:characteristic-eq-is-poly}).
Kiedy to ustalimy,
	wyznaczyć możemy $\ker(A - \lambda I)$.
Każdy wektor w tej przestrzeni liniowej będzie wektorem własnym $A$.
Wielomian,
	który pojawia się w tym równaniu, ma swoją specjalną nazwę.

\begin{definition}
	\textit{Wielomianem charakterystycznym} macierzy $A \in \mathbb{C}^{n \times n}$
		nazywamy wielomian
	\begin{equation}
		\chi_A(t) = \det(t I - A).
	\end{equation}
\end{definition}

\begin{definition}
	Krotnością algebraiczną wartości własnej $\lambda$ macierzy $A \in \mathbb{C}^{n \times n}$
		nazywamy krotność pierwiastka $\lambda$ wielomianu $\chi_A$.
\end{definition}

\begin{example}
	Niech
	\begin{equation}
		A = \left[\begin{matrix}
		\end{matrix}\right].
	\end{equation}
	Wyznaczymy wartości własne macierzy $A$ oraz ich krotności.
\end{example}

% TODO: 0 v \textbf{0}


\begin{definition}
	Macierz $A \in \mathbb{C}^{n \times n}$ nazywamy \textit{diagonalizowalną},
		jeśli istnieje macierz odwracalna $X \in \mathbb{C}^{n \times n}$
		i macierz diagonalna $D \in \mathbb{C}^{n \times n}$,
		taka, że
	\begin{equation}
		A = X D X^{-1}.
	\end{equation}
\end{definition}
Macierz $X$ w powyższym równaniu traktujemy jako macierz zmiany bazy,
	a o macierzy $D$ mówimy wtedy jako o postaci diagonalnej macierzy $A$.
Dla danej macierzy $A$ istnieć może wiele różnych postaci diagonalnych!
Nie są one unikalne.

Zauważmy,
	że macierz jest diagonalizowalna wtedy i tylko wtedy,
	gdy ma $n$ wektorów własnych,
	których zbiór jest liniowo niezależny.
% TODO: dlaczego?


\begin{example}
	Macierz
	\begin{equation}
		A = \left[\begin{matrix}
			1 & 1 & 0 \\
			1 & 1 & 0 \\
			0 & 0 & 2
		\end{matrix}\right]
	\end{equation}
	jest diagonalizowalna.
	Wyznaczmy jej wartości własne.
\end{example}

\begin{example}
	Macierz
	\begin{equation}
		A = \left[\begin{matrix}
			2 & 1 & 0 \\
			0 & 2 & 1 \\
			0 & 0 & 2
		\end{matrix}\right]
	\end{equation}
	nie jest diagonalizowalna.
	
\end{example}

% TODO: w postaci diagonalnej na diagonali są wektory własne

\subsection{Zadania.}

\begin{exercise}\label{ex:characteristic-eq-is-poly}
	Niech $A \in \mathbb{C}^{n \times n}$ będzie macierzą.
	Udowodnij,
		że $\det(\lambda I - A)$ jest wielomianem z $\lambda$ jako zmienną.
\end{exercise}

\begin{exercise}
	Niech $A \in \mathbb{C}^{n \times n}$ będzie macierzą.
	Udowodnij,
		że $\chi_A(t)$ jest wielomianem \textit{monicznym},
		czyli że jego wiodący współczynnik jest równy $1$.
\end{exercise}
